
\subsection{State Transition Functions}\label{sec:stf}

State Transition Functions (STF) define the transition of states from a pre-event to a post-event state when a certain event $e_t^{k}$ applies. Thereby, the pre-event and post-event times are indexed with $t^-$ and $t^+$, respectively. 

These functions are specific to a certain event and contract. STFs are written according to the following pattern \stf{[event type]}{[contract type]} where [event type] and [contract type] refer to the respective event type and contract to which the STF belongs.

\begin{example}[State Transition Functions]
The STF for an IP event and PAM contract is written as \stf{IP}{PAM} and maps e.g. state \textit{Accrued Interest} from a pre-event state $\svar{Ipac}{t^-}$ to post-event state $\svar{Ipac}{t^+}$.
\end{example}


\subsection{Payoff Functions}\label{sec:pof}

Payoff Functions (POF) define how the cash flow $c\in\Real$ for a certain event $e_t^k$ is being derived from current states and from the contract terms. If necessary, the resulting cash flow can be indexed with the event time $c_t$. These functions are specific to a certain event and contract. POFs are written according to the following pattern \pof{[event type]}{[contract type]} where [event type] and [contract type] refer to the respective event and contract to which the STF belongs.

\begin{example}[Payoff Functions]
The POF for an IP event $e_t^{IP}$ and PAM contract is written as \pof{IP}{PAM} with $\fev{e_t^{IP}}=\pof{IP}{PAM}$.
\end{example}
